\note{sa page 4}

$u_{2n-i}$ est \struck{sous} \ill{} la conjonction des suivantes :

\[ (3.1)_{i} \qquad \bigl( u \in \mathfrak{z}\mathcal{E}^{0}, \ \mathcal{E}^{2n-i}_{i}\, u = 0 \bigr) \implies \pi_{i}\, u = 0 \]
\[ (3.2)_{i} \qquad \bigl( u \in \mathfrak{z}\mathcal{E}^{0}, \ u\, \mathcal{E}^{2n-i}_{i} = 0 \bigr) \implies \pi_{2n-i}\, u = 0 \]

En effet, la relation $(3.1)$ implique que \struck{$\pi_{i} \in \mathfrak{z}\mathcal{E}^{0}$
est} l'annulateur de $\pi_{i}$ dans $\mathfrak{z}\mathcal{E}^{0}$ est l'idéal
des $u \in \mathfrak{z}\mathcal{E}^{0}$ tels que $\mathcal{E}^{2n-i}_{i} u = 0$ ;
et dans un anneau commutatif, un idempotent est connu quand on connaît son
annulateur. De même pour $\pi_{2n-i}$, son annulateur dans
$\mathfrak{z}\mathcal{E}^{0}$ est formé des $u \in \mathfrak{z}\mathcal{E}^{0}$
tels que $u\, \mathcal{E}^{2n-i}_{i} = 0$.

En même temps, les conditions $(3.1)_{i}$, $(3.2)_{i}$, pour tout $i$, donnent
une \struck{condition} possibilité de déterminer de proche en proche les
$\pi_{i}$.

4. Lorsque $\mathcal{E}^{*}$ est comme dans 1. \add{ou v. (2.2.)}, une condition
suffisante pour avoir $(3.1)_{i}$ est la \struck{\ill{}} suivante :

\[ (4.1)_{i} \qquad \text{l'intersection des noyaux des } u \colon M^{i} \to M^{2n-i}, \ u \in \mathcal{E}^{2n-i}_{i}, \ \text{est réduite à } 0 . \]
\[ (4.2)_{i} \qquad \text{la somme des images des } u \colon M^{i} \to M^{2n-i}, \ u \in \mathcal{E}^{2n-i}_{i}, \ \text{est égale à } M^{2n-i} . \]

Ces deux conditions s'expriment \ill{}

\[ (4.3)_{i} \qquad \text{il existe } u \in \mathcal{E}^{2n-i}_{i} \ \text{qui soit un isomorphisme} \ \ill{}, \ \text{ou, de} \]

